\chapter*{Introduction générale}
\addcontentsline{toc}{chapter}{Introduction générale}
\markboth{Introduction générale}{Introduction générale}
\label{ch:intro}

\usection{Contexte et motivation}
L'algorithme de Shor~\cite{shor} rend les signatures fondées sur la factorisation et le logarithme discret vulnérables à un ordinateur quantique de grande taille. Le NIST a publié en 2024 les premières normes post-quantiques (FIPS~203, 204 et 205~\cite{nistfips}) et poursuit un appel à signatures additionnelles~\cite{nistadd}, dans lequel figurent des schémas multivariés. Parmi eux, MAYO~\cite{mayo} construit une signature de type Oil-and-Vinegar (OV)~\cite{patarin,kipnisshamir,uov} à clé publique compacte : l'application publique est une application quadratique d'un petit espace d'huile, « fouettée » (\emph{whipped}) pour produire une application plus grande. L'intérêt pratique du schéma est la petite taille des clés et des signatures ; son coût de calcul est dominé par de l'arithmétique sur $\mathbb{F}_{16}$.

Les schémas multivariés ont connu des cryptanalyses marquantes~\cite{beullensrainbow}, ce qui impose une lecture prudente de toute affirmation de sécurité ; ce mémoire ne traite pas de sécurité algébrique. Il traite de \emph{performance d'implémentation} et de \emph{propriétés d'implémentation} (correction, accès mémoire, temps constant) d'une optimisation locale du code C officiel.

\paragraph{Parcours du travail et recentrage.} Le travail a débuté par une implémentation autonome (arithmétique $\mathbb{F}_{16}$ à tables, étude d'une fuite par canal auxiliaire de cache sur les tables, élimination de Gauss--Jordan à temps constant), documentée dans des chapitres antérieurs et motivée par les attaques temporelles de cache~\cite{kocher,bernstein}. À la demande de l'encadrant, le travail a ensuite été recentré sur l'implémentation C \emph{officielle} (MAYO-C) pour mesurer des gains qui soient comparables à l'existant. Ce mémoire porte exclusivement sur cette seconde phase. L'étude préliminaire n'est pas reprise, sauf quand un résultat négatif de cette phase éclaire une décision (section~\ref{sec:discussion}).

\usection{Problématique}
Dans l'implémentation AVX2 officielle, un $m$-vecteur de $\mathbb{F}_{16}$ est stocké à deux éléments par octet ($m/2$ octets) et multiplié par des scalaires à l'aide de \vpshufb{} : une instruction effectue 32 recherches dans une table de 16 octets et fournit, pour chaque élément, les produits par \emph{deux} scalaires à la fois. Un registre de 256 bits couvre donc 64 éléments. Quand $m$ n'est pas un multiple de 64, le code officiel couvre la fin du vecteur par un chargement de 32 octets reculé qui \emph{recouvre} le précédent. Ce recouvrement est documenté dans l'article d'optimisation~\cite{nibbling} pour $m=96$ (MAYO-3) : il « duplique » des opérations.

Pour les paramètres considérés ici, la duplication est importante : 48 éléments recalculés sur 128 pour MAYO-1 ($m=80$, soit 37,5\,\%) et 50 sur 192 pour MAYO-5 ($m=142$, soit 26\,\%). La question centrale est :

\begin{encadre}
\textbf{Question centrale.} La duplication de calcul due au recouvrement en fin de $m$-vecteur peut-elle être réduite par un traitement dédié de la queue, quelle en est la limite exacte en nombre de \vpshufb{}, et ce gain se traduit-il, de façon robuste et démontrable, en gain sur le temps de signature complet ?
\end{encadre}

Cette formulation sépare trois choses souvent confondues : (i) une \emph{borne de comptage} (mathématique, dans un modèle précis), (ii) sa \emph{réalisation} par un algorithme (à montrer et à tester) et (iii) l'\emph{effet sur le temps} (expérimental, dépendant du compilateur et de la machine).

\usection{Objectifs de l'étude}
\begin{description}[leftmargin=2.2cm,labelwidth=2cm]
\item[QR1] Quel est le nombre minimal de \vpshufb{} par paire de scalaires pour un $m$-vecteur, dans la classe des traitements à tables de 16 octets, et pour quels $m$ le code officiel est-il déjà optimal ?
\item[QR2] Existe-t-il une réalisation sûre (sans lecture hors de l'entrée, sans dépendance au secret dans les adresses) de cette borne pour MAYO-1 et MAYO-5 ?
\item[QR3] Quel gain mesure-t-on sur les noyaux \code{P1_times_Vt} et \code{P1t_times_Vt}, puis sur \code{mayo_sign}, sur une machine physique, avec un protocole dont les limites sont quantifiées ?
\item[QR4] Si le gain en contexte diffère du gain isolé, que montre l'analyse du code généré, et qu'est-ce qui reste à démontrer ?
\end{description}

Les objectifs en découlent : \textbf{O1} formaliser le comptage (borne, réalisation, limite) ; \textbf{O2} documenter la conception d'A2 et les propriétés dont dépend sa correction ; \textbf{O3} mesurer avec un protocole reproductible ; \textbf{O4} caractériser l'écart entre gain isolé et gain en contexte ; \textbf{O5} situer le travail par rapport aux implémentations existantes sans revendiquer plus que ce que les sources consultées permettent.

Conformément à la règle de rédaction adoptée, le tableau~\ref{tab:contrib} sépare ce qui est acquis de ce qui reste ouvert.

\begin{table}[ht]\centering\small
\caption{Statut des contributions.}\label{tab:contrib}
\begin{tabularx}{\linewidth}{@{}lXl@{}}
\toprule
 & Contribution & Statut\\\midrule
C1 & Comptage : borne $\lceil m/32\rceil$, schéma officiel $2\lceil m/64\rceil$, écart $\Delta(m)$ (prop.~\ref{prop:borne}--\ref{prop:real}) & établi (preuve)\\
C2 & Réalisation A2 (queue $\le 16$ éléments) pour MAYO-1 et MAYO-5, variante SAFE sans lecture hors de P1 & établi (code, KAT, tests)\\
C3 & Gains sur noyaux isolés (12--29\,\%) et sur \code{mayo_sign} (MAYO-1 : 10,3\,\% ; MAYO-5 : 1,3\,\%) sur une machine, un compilateur & établi, avec limites\\
C4 & Localisation de la perte de gain de MAYO-5 dans \code{m_calculate} en contexte de signature & établi (mesure directe)\\
C5 & Écritures supplémentaires asymétriques dans le code généré par GCC~16.2.0, cohérentes avec C4 ; identifiées (E1b) à des stockages de \code{temp[]}/\code{tt[]} dans la boucle ; SRA écartée & fait statique établi ; causalité sur le temps \emph{non} établie\\
C6 & Dépôt public de scripts, données brutes et empreintes, reproduisant les analyses (9/9) & établi\\\midrule
Q1 & Pourquoi GCC ne sort pas ces stockages de la boucle sans toutes les sondes (passe responsable) & \emph{ouvert} (dumps \code{lim}/\code{pre} non produits)\\
Q2 & Lien causal écritures $\to$ perte de gain, avec un modèle qui n'en surestime pas l'effet & \emph{ouvert}\\
Q3 & Variante pour queues de 17 à 32 éléments ; autres compilateurs, autres microarchitectures & \emph{ouvert} (non implémenté)\\
Q4 & Temps constant au sens fort (preuve formelle) & \emph{non fourni}\\\bottomrule
\end{tabularx}
\end{table}

\usection{Méthodologie de recherche}
La démarche est expérimentale et déductive. Le comptage est établi par preuve dans un modèle défini ; la réalisation est justifiée par le code et testée ; les gains sont mesurés sur machine physique avec un protocole apparié, dont l'unité statistique est le processus ; le code généré par le compilateur est analysé statiquement pour interpréter les écarts. Chaque résultat est étiqueté selon son statut (fait, inférence, hypothèse), et les scripts, données brutes et empreintes sont déposés pour permettre la reproduction.

Le travail est limité aux noyaux \code{P1_times_Vt} et \code{P1t_times_Vt} pour MAYO-1 et MAYO-5. MAYO-2 ($m=64$) et MAYO-3 ($m=108$) servent de témoins : pour eux, aucun gain n'est possible dans la classe étudiée (proposition~\ref{prop:borne}). Aucune autre optimisation n'est proposée, les noyaux optimisés ne sont pas modifiés pendant l'analyse, et la décision d'intégration dans MAYO-C est volontairement différée. L'absence d'une technique dans les sources consultées ne prouve pas qu'elle n'existe pas ailleurs : aucune nouveauté absolue n'est revendiquée.

Les énoncés portant sur des résultats sont étiquetés : \F{} \emph{fait} (lu dans le code généré ou dans des mesures déposées), \I{} \emph{inférence} tirée des faits, \Hy{} \emph{hypothèse} non testée. Les propriétés de sûreté du section~\ref{sec:surete} utilisent \Pc{} (justifiée par le code) et \Pe{} (vérifiée expérimentalement) ; une propriété \Pe{} seule n'est pas une démonstration. Les durées sont en ticks du compteur TSC (rdtscp), non en cycles cœur.

\usection{Intérêt de l'étude}
Pour les mainteneurs et les utilisateurs de MAYO-C, l'étude fournit une borne précise sur ce qu'une optimisation du recouvrement peut gagner en nombre de \vpshufb{}, une réalisation vérifiée pour deux jeux de paramètres, et surtout une mise en garde méthodologique : un gain mesuré sur un noyau isolé ne se transfère pas automatiquement à la signature, car le code généré par le compilateur dépend du contexte. Pour la méthode, le mémoire propose un protocole de mesure et d'analyse statique qui sépare les faits des hypothèses et qui est reproductible.

\usection{Structure du mémoire}
Le mémoire comprend trois chapitres, encadrés par cette introduction et par une conclusion. Le chapitre~\ref{ch:etat} présente MAYO, l'arithmétique par \vpshufb{} et les implémentations existantes, et situe A2. Le chapitre~\ref{ch:methodo} formalise le comptage et décrit A2 (section~\ref{sec:a2}), établit ses propriétés de correction et de sûreté (section~\ref{sec:surete}) et détaille le protocole expérimental (section~\ref{sec:protocole}). Le chapitre~\ref{ch:resultats} présente les résultats et les discute en séparant faits, inférences et hypothèses. La conclusion récapitule contributions, limites et perspectives.
